\chapter{Conclusiones y trabajos futuros}
\label{ch:conclusiones}

\section{Introducción del capítulo}
Toda comparación técnica necesita cerrarse con una lectura de alcance y con una agenda clara
de continuidad. En el caso de este libro, la discusión final no busca únicamente afirmar que un
algoritmo fue mejor que otro, sino precisar qué se aprendió sobre la sintonización PID en
microredes aisladas y cuáles son los siguientes pasos razonables para convertir ese conocimiento
en una herramienta de diseño más madura.

\section{Objetivo del capítulo}
Este capítulo sintetiza los hallazgos centrales del estudio, explicita las contribuciones más
importantes y delimita los aspectos que aún requieren mayor validación. También propone
líneas de trabajo futuro orientadas a pruebas experimentales, comparación con algoritmos
adicionales e integración con plataformas de implementación más cercanas a la industria.

\section{Conclusiones generales}
El hallazgo principal es que la sintonización PID basada en optimización metaheurística ofrece
una mejora consistente frente a los ajustes clásicos cuando la planta presenta no linealidades y
perturbaciones severas. El material base centrado en PSO muestra reducciones claras en ITAE,
tiempo de asentamiento y oscilación de frecuencia, lo que respalda la hipótesis de que una
búsqueda guiada por desempeño supera a una parametrización obtenida solo a partir de
ganancia crítica y periodo crítico.

\section{Aportes del libro}
Entre los aportes más relevantes destacan tres. En primer lugar, se organiza una ruta de
trabajo completa que une fundamentos del PID, criterios de evaluación y formulación del
problema como optimización. En segundo lugar, se adapta ese marco al caso de una microred
hidro-solar andina, aportando una lectura aplicada y no meramente algorítmica. En tercer lugar,
se deja preparada una estructura comparativa capaz de incorporar no solo PSO, sino también
GWO, Eagle Strategy y ABC dentro de un mismo esquema metodológico.

\section{Limitaciones}
El alcance de esta obra es principalmente metodológico y de simulación. Los resultados son
valiosos para entender tendencias dinámicas, pero todavía dependen de la fidelidad del modelo,
de la calidad de los parámetros de planta y del entorno numérico empleado. En consecuencia,
no deben interpretarse como una validación definitiva para implementación industrial directa,
sino como una base robusta para avanzar hacia pruebas más cercanas a operación real.

\section{Líneas futuras de investigación}
Las extensiones naturales del estudio incluyen validación Hardware-in-the-Loop, pruebas con
PLC o controladores industriales, análisis de estabilidad de pequeña señal y campañas
comparativas con funciones objetivo multiobjetivo. También resulta pertinente explorar
hibridaciones entre algoritmos, variantes adaptativas de sintonización y estrategias que integren
restricciones energéticas más explícitas, por ejemplo bajo escenarios prolongados de estiaje o
alta penetración solar.

\section{Cierre del capítulo}
La principal conclusión de fondo es que la modernización del control no requiere necesariamente
reemplazar la estructura PID, sino refinar la manera en que se ajusta. Cuando esa sintonización
se apoya en métricas pertinentes, modelos consistentes y métodos de búsqueda adecuados, es
posible obtener mejoras sustanciales con una inversión razonable y con impacto directo sobre la
confiabilidad de la electrificación rural.
